\subsection{04.05.2026 (Freistunde)}

\subsubsection{Ziel}

Ziel des Tages war es, die Datenbankstruktur
ordnungsgemäß umzusetzen und die Speicherung
der Messdaten vorzubereiten.

\subsubsection{Durchgeführte Arbeiten}

Zunächst wurde die Datei \texttt{datenbank.py}
erstellt.
Diese dient als allgemeine Schnittstelle für die
Verbindung zur Datenbank sowie für grundlegende
Datenbankoperationen.

Zusätzlich wurde die Datei
\texttt{taupunkt\_db.py} erstellt.
Darin wurden spezielle Methoden für die
Datenbankstruktur des Projekts implementiert.

Außerdem wurde eine Tabelle vorbereitet,
in der sämtliche Messdaten gespeichert werden
sollen.

Die Tabelle enthält folgende Werte:

\begin{itemize}
    \item Temperatur innen
    \item Temperatur außen
    \item Luftfeuchtigkeit innen
    \item Luftfeuchtigkeit außen
    \item Taupunkttemperatur innen
    \item Taupunkttemperatur außen
    \item Taupunktdifferenz
    \item Zeitpunkt
\end{itemize}

Dabei wurde darauf geachtet, die Datenbankstruktur
möglichst übersichtlich und erweiterbar aufzubauen.

\subsubsection{Erkenntnisse}

Es wurde festgestellt, dass eine Trennung zwischen
allgemeiner Datenbanklogik und projektspezifischen
Datenbankmethoden die Übersichtlichkeit und
Wartbarkeit des Projekts verbessert.

Zusätzlich wurde deutlich, dass eine saubere
Tabellenstruktur wichtig ist, damit Messdaten
später einfach gespeichert und ausgewertet
werden können.

\subsubsection{Nächste Schritte}

Am nächsten Tag soll getestet werden,
ob die Datenbank korrekt mit den
Messdaten funktioniert und die Werte
erfolgreich gespeichert werden können.